\documentclass[12pt]{article}

\usepackage{fullpage}
\usepackage{graphicx}
\usepackage{amssymb}
\usepackage{amsmath}
\usepackage[none]{hyphenat}
\usepackage{parskip}
\usepackage[spanish]{babel}
\usepackage[utf8]{inputenc}
\usepackage{hyperref}
\usepackage{fancyhdr}
\usepackage{tasks}
\usepackage{mdframed}
\usepackage{xcolor}
\usepackage{pgfplots}
\usepackage[makeroom]{cancel}
\usepackage{multicol}
\usepackage[shortlabels]{enumitem}
\usepackage{stackrel}
\usepackage{tkz-tab}
\usepackage{xpatch}
\xpatchcmd{\tkzTabLine}{$0$}{$\bullet$}{}{}

\setlength{\headheight}{10pt}
\setlength{\headsep}{10pt}
\pagestyle{fancy}
\rhead{\ayudantia \ - \alumno}
\tikzset{t style/.style={style=solid}}

\newcommand*{\mybox}[2]{\colorbox{#1!30}{\parbox{.98\linewidth}{#2}}}

\newenvironment{solucion}
{\begin{mdframed}[backgroundcolor=black!10]
		{\bf Solución:}\\
	}
	{
	\end{mdframed}
}

\newenvironment{alternativas}[1]
{\begin{multicols}{#1}
		\begin{enumerate}[a)]
		}
		{
		\end{enumerate}
	\end{multicols}
}

\newenvironment{preguntas}
{\begin{enumerate}\itemsep12pt
	}
	{
	\end{enumerate}
}

\newcommand{\ayudantia}{{\sc Ayudantía 12.5}}
\newcommand{\tituloayu}{Compilado I3}
\newcommand{\fecha}{7 de junio de 2019}
\newcommand{\sigla}{MAT1203}
\newcommand{\nombre}{Álgebra Lineal}
\newcommand{\profesor}{Camilo Perez}
\newcommand{\ano}{2019}
\newcommand{\semestre}{1}
\newcommand{\mail}{mat1203@ifcastaneda.cl}
\newcommand{\alumno}{Ignacio Castañeda - \mail}

\newcommand{\ev}{\Big|}
\newcommand{\ra}{\rightarrow}
\newcommand{\lra}{\leftrightarrow}
\newcommand{\N}{\mathbb{N}}
\newcommand{\R}{\mathbb{R}}
\newcommand{\Exp}[1]{\mathcal{E}_{#1}}
\newcommand{\List}[1]{\mathcal{L}_{#1}}
\newcommand{\EN}{\Exp{\N}}
\newcommand{\LN}{\List{\N}}
\newcommand{\comment}[1]{}
\newcommand{\lb}{\\~\\}
\newcommand{\eop}{_{\square}}
\newcommand{\hsig}{\hat{\sigma}}
\newcommand{\widesim}[2][1.5]{
	\mathrel{\overset{#2}{\scalebox{#1}[1]{$\sim$}}}
}
\newcommand{\wsim}{\widesim{}}
\newcommand{\lh}{\stackrel{L'H}{=}}

\begin{document}
\thispagestyle{empty}

\begin{minipage}{2cm}
	\includegraphics[width=2cm]{../../../../img/logo.pdf}
	\vspace{0.5cm}
\end{minipage}
\begin{minipage}{\linewidth}
	\begin{tabular}{lrl}
		{\scriptsize\sc Pontificia Universidad Catolica de Chile} & \hspace*{0.7in}Curso: &
		\sigla  - \nombre\\
		{\sc Facultad de Matemáticas}&
		Profesor: & \profesor \\
		{\sc Semestre \ano-\semestre} & Ayudante: & {Ignacio Castañeda}\\
		& {Mail:} & \texttt{\mail}
	\end{tabular}
\end{minipage}

\vspace{-10mm}
\begin{center}
	{\LARGE\bf \ayudantia}\\
	\vspace{0.1cm}
	{\tituloayu}\\
	\vspace{0.1cm}
	\fecha\\
	\vspace{0.4cm}
\end{center}

\begin{preguntas}
\item Sea 
	$$U=\{p(x) \in \mathbb{P}_2 : p(1) + p(0) = p(-1)\}$$
	Determinar una base de $U$.
\begin{solucion}
			Sea $p(x) \in \mathbb{P}_2$ un polinomio de la forma $p(x) = a + bx + cx^2$\\
			\\
			Para que pertenezca a $U$, debe cumplirse que
			$$p(1) + p(0) = p(-1)$$
			$$a+b+c+a=a-b+c$$
			$$a = -2b$$
			Luego,
			$$\begin{array}{rcl}
			a & = & -2b\\
			b & = & b\\
			c & = & c
			\end{array}$$
			Entonces, podemos reescribir el polinomio como
			$$p(x) = -2b + bx + cx^2$$
			$$p(x) = b(x-2) + c(x^2)$$
			Finalmente,
			$$U = Gen\{x-2, x^2\}$$
\end{solucion}
\item Use vectores de coordenadas para probar la independencia lineal del conjunto de polinomios $\{1-2t^2-t^3,t+2t^3,1+t-2t^2\}$.
\begin{solucion}
La base que utilizaremos será
		$$B = \{1,t,t^2,t^3\}$$
		Luego, podemos reescribir cada elemento del conjunto de la siguiente forma
		$$\begin{array}{lcl}
		1-2t-t^3 & \ra & (1,\ \ 0,-2,-1)\\
		t+2t^3 & \ra & (0,\ \ 1,\ \ 0,\ \ 2)\\
		1+t-2t^3 & \ra & (1,\ \ 1,-2,\ \ 0)
		\end{array}$$
		Ahora, para probar que los elementos del conjunto son $L.I.$, basta con probar que los vectores de coordenadas lo son. Para esto, ponemos los vectores en una matriz y pivoteamos:
		$$\begin{bmatrix}
		1 & 0 & 1\\
		0 & 1 & 1\\
		-2 & 0 & -2\\
		-1 & 2 & 0
		\end{bmatrix} \sim
		\begin{bmatrix}
		1 & 0 & 1\\
		0 & 1 & 1\\
		0 & 0 & 1\\
		0 & 0 & 0
		\end{bmatrix}$$
		Con lo que concluimos que son $L.I.$
		$$\blacksquare$$
\end{solucion}
\item Sean $p_1(t) = 1-t^2, p_2(t) = 1+t, p_3(t) = 1+t+t^2$. Se sabe que $\{p_1(t), p_2(t), p_3(t)\}$ es una base de $\mathbb{P}_2$.
\begin{enumerate}[a)]
\item Exprese los polinomios $f(t) = 3-5t + 2t^2$ y $g(t) = 1-3t$ como combinaciones lineales de $\{p_1(t), p_2(t), p_3(t)\}$.
\item Use los vectores de coordenadas encontrados en la parte anterior para determinar si el conjunto $\{p_1(t), f(t), g(t)\}$ es L.I. o L.D.
\end{enumerate}
\begin{solucion}

\begin{enumerate}[a)]
\item Exprese los polinomios $f(t) = 3-5t + 2t^2$ y $g(t) = 1-3t$ como combinaciones lineales de $\{p_1(t), p_2(t), p_3(t)\}$.\\
			\\
			Sea una base 
			$$B = \{p_1(t), p_2(t), p_3(t)\} = $$, lo que estamos buscando son las coordenadas de $f$ y $g$ en la base $B$.\\
			\\
			En primer lugar, veamos $f(t)$. Debemos buscar $a$, $b$, $c$, tal que
			$$f(t) = 3-5t + 2t^2 = ap_1(t) + bp_2(t) + cp_3(t)$$
			Reemplazando y reagrupando,
			$$3-5t + 2t^2 = a(1-t^2) + b(1+t) + c(1+t+t^2)$$
			$$3-5t + 2t^2 = (a+b+c) + (b+c)t + (c-a)t^2$$
			Por lo que
			$$\begin{array}{rl}
			a + b +c & = 3\\
			b + c & =-5\\
			-a+c & =2
			\end{array} \Longrightarrow
			\begin{array}{rl}
			a & = 8\\
			b & =-15\\
			c & = 10
			\end{array}$$
			Luego, las coordenadas de $f(t)$ en $B$ son $(8,-15,10)$.\\
			\\
			Veamos ahora que ocurre con $g(t)$. Nuevamente, buscamos $a$, $b$, $c$, tal que
			$$g(t) = 1-3t =  ap_1(t) + bp_2(t) + cp_3(t)$$
			Esto es,
			$$1-3t = a(1-t^2) + b(1+t) + c(1+t+t^2)$$
			$$1-3t = (a+b+c) + (b+c)t + (c-a)t^2$$
			Lo que se nos lleva al sistema
			$$\begin{array}{rl}
			a + b +c & = 1\\
			b + c & =-3\\
			-a+c & =0
			\end{array} \Longrightarrow
			\begin{array}{rl}
			a & = 4\\
			b & =-7\\
			c & = 4
			\end{array}$$
			Con lo que las coordenadas de $g(t)$ en $B$ son $(4,-7,4)$.
\item Use los vectores de coordenadas encontrados en la parte anterior para determinar si el conjunto $\{p_1(t), f(t), g(t)\}$ es L.I. o L.D.\\
			\\
			Para determinar esto, podemos usar las coordenadas de estos elementos en cualquier base. Por conveniencia, usaremos la base $B$. Aquí, los vectores de coordenadas son
			$$p_1(t) = (1,0,0)$$
			$$f(t) = (8,-15,10)$$
			$$g(t) = (4,-7,4)$$
			Armamos una matriz con ellos y pivoteamos,
			$$\begin{bmatrix}
			1 & 0 & 0 \\
			8 & -15 & 10 \\
			4 & -7 & 4
			\end{bmatrix}
			\sim
			\begin{bmatrix}
			1 & 0 & 0 \\
			0 & 1 & 0 \\
			0 & 0 & 1
			\end{bmatrix}$$
			Por lo que sus el conjunto $\{p_1(t), f(t), g(t)\}$ es $L.I.$
\end{enumerate}
\end{solucion}
\item Sea
$$T\left(\left[ \begin{array}{c}
a\\ b\\ c \end{array} \right] \right) = (a+b+c) + (a-b+2c)x +
(3b-c)x^2
$$
Determine una base para $Nul(T) $ y para $Im (T) $ y las dimensiones de estos subespacios.
\begin{solucion}
Notemos que lo que estamos generando son polinomios de grado menor o igual a 2. La base canónica para estos polinomios es $B=\{1, x, x^2\}$
En primer lugar, obtengamos la matriz asociada a la transformación lineal respecto a las bases canónicas, esto es,
$$A = \begin{bmatrix}
1 & 1 & 1 \\
1 & -1 & 2 \\
0 & 3 & -1
\end{bmatrix}$$
Recordemos que $Nul(T)$ corresponde a la solución del sistema homogeneo, es decir $Ax = 0$, esto es,
$$\begin{bmatrix}
1 & 1 & 1 \\
1 & -1 & 2 \\
0 & 3 & -1
\end{bmatrix} \sim 
\begin{bmatrix}
1 & 0 & 0 \\
0 & 1 & 0 \\
0 & 0 & 1
\end{bmatrix}$$
Por lo que la solución del sistema homogeneo es $Gen\left\{\begin{pmatrix}0 \\ 0 \\ 0 \end{pmatrix}\right\}$. Recordemos que este vector es un vector de coordenadas, sin embargo, al ser el vector nulo, este generado equivale a $Gen\{0\}$.\\

Veamos ahora el $Im (T)$, que corresponde a las bases $LI$. Esto es equivalente a ver las columnas $LI$ de $A$, que ya sabemos que son todas. Por lo tanto, 
$$Im(A) = Gen\left\{\begin{pmatrix}1 \\ 0 \\ 0 \end{pmatrix},\begin{pmatrix}0 \\ 1 \\ 0 \end{pmatrix}, \begin{pmatrix}1 \\ 0 \\ 0 \end{pmatrix}\right\}$$
Luego, traduciendo estos vectores coordenadas, concluimos que
$$Im(T) = Gen\{1, x, x^2\}$$
Finalmente, las dimensiones de $Nul(T)$ y $Im(T)$ son 0 y 3, respectivamente.
\end{solucion}
\item Sean $B_1 = \{1+x,1-x,1+x^2\}$ y $B_2$ bases de $P_2(\R)$ tales que para todo $p \in P_2(\R)$
	$$ \begin{bmatrix}
	1 & 0 & -1\\
	1 & 1 & 0 \\
	0 & 1 & 2
	\end{bmatrix} [p]_{B_1} = [p]_{B_2}$$
	Determine los polinomios que forman la base $B_2$.
\begin{solucion}
Notemos que la matriz que nos dan corresponde a la matriz cambio de base de $B_1$ a $B_2$. Recordemos que una matriz cambio de base, de $B_1$ a $B_2$, corresponde a la matriz formada por las coordenadas de los elementos de $B_1$ en la base $B_2$.
		
		En base a lo anterior, dado que tenemos los vectores de $B_1$, lo que necesitamos es justamente lo contrario, es decir, la matriz cambio de base de $B_2$ a $B_1$. De esta manera, tendriamos las coordenadas de los elementos de $B_2$ en la base $B_1$, por lo que podriamos calcularlos con facilidad, ya que tenemos estos últimos.
		
		Para encontrar esta matriz, basta con obtener la inversa de la matriz cambio de base dada, esto es
		$$ \left[\begin{array}{ccc|ccc}
		1 & 0 & -1 & 1 & 0 & 0\\
		1 & 1 & 0 & 0 & 1 & 0\\
		0 & 1 & 2 & 0 & 0 & 1
		\end{array}\right] \sim 
		\left[\begin{array}{ccc|ccc}
		1 & 0 & 0 & 2 & -1 & 1\\
		0 & 1 & 0 & -2 & 2 & -1\\
		0 & 0 & 1 & 1 & -1 & 1
		\end{array}\right]$$
		con lo que
		$$A_{B_2}^{B_1} = \begin{bmatrix}
		2 & -1 & 1\\
		-2 & 2 & -1\\
		1 & -1 & 1
		\end{bmatrix}$$
		Sea 
		$$B_2 = \{q_1(x), q_2(x), q_3(x)\}$$
		Las columnas de $A_{B_2}^{B_1}$ corresponden a los vectores de coordenada de $q_1$, $q_2$ y $q_3$, respectivamente, en la base $B_1$.
		
		Luego,
		$$q_1(x) = 2(1+x) - 2(1-x) + (1+x^2) = 1+4x+x^2$$
		$$q_2(x) = -(1+x) + 2(1-x) - (1+x^2) = -3x-x^2$$
		$$q_3(x) = (1+x) - (1-x) + (1+x^2) = 1+2x+x^2$$
		Finalmente,
		$$B_2 = \{1+4x+x^2,-3x-x^2,1+2x+x^2\}$$
\end{solucion}
\item Sean $B$ y $C$ bases de un espacio vectorial $V$ y $P = \begin{bmatrix}4 &-1 \\ 6 & -1\end{bmatrix}$ la matriz de cambio de coordenadas tal que $[v]_C = P[v]_B\ \forall v \in V$
\begin{enumerate}[a)]
\item Demuestre que el conjunto $W = \{v \in V: [v]_C = 2[v]_B\}$ es un subespacio de $V$.
\item Si $B=\{v_1,v_2\}$ determine una base para $W$ en términos de la base de $B$.
\end{enumerate}
\begin{solucion}

\begin{enumerate}[a)]
\item Demuestre que el conjunto $W = \{v \in V: [v]_C = 2[v]_B\}$ es subespacio de $V$.
			
			\begin{enumerate}
				\item No vacio $(0\in W)$
				$$\begin{array}{cc}
				[\vec{0}]_C = \vec{0}\\
				2 \cdot [\vec{0}]_B = 2 \cdot \vec{0} = 0
				\end{array} \Longrightarrow \vec{0} = \vec{0} \ra \vec{0} \in W$$	
				\item Suma y multiplicación (a la vez)	
				
				
				Sea $u, v \in W$ y $\alpha, \beta \in \R$, debemos demostrar que $\alpha u + \beta v \in W$
				$$\begin{array}{rcl}
				[\alpha u + \beta v]_C & = & 2 [\alpha u + \beta v]_B \\
				& = & 2 [\alpha u + \beta v]_B \\
				& = & 2 [\alpha u]_B + 2[\beta v]_B \\
				& = & \alpha 2[u]_B + \beta 2[v]_B \\
				& = & \alpha [u]_C + \beta [v]_C \\
				& = & [\alpha u]_C + [\beta v]_C \\
				{[}\alpha u + \beta v]_C & = & [\alpha u +\beta v]_C
				\end{array}
				$$
			\end{enumerate}
			$$\blacksquare$$
\item Si $B=\{v_1,v_2\}$ determine una base para $W$ en términos de la base de $B$.
			
			Sea $B = \{v_1, v_2\}$ la base dada, busquemos una base $B_W$ de $W$\\
			\\
			Sea $v \in W$,
			$$v = x_1v_1 + x_2v_2 \ra x = \begin{pmatrix} x_1 \\ x_2 \end{pmatrix} = [v]_B$$
			Además,
			$$\begin{array}{cclcl}
			[v]_C & = & 2[v]_B & = & 2x \\
			{[}v]_C & = & P[v]_B & = & Px
			\end{array} \Longrightarrow
			2x = Px$$
			Luego,
			$$2x = Px$$
			$$Px - 2x = 0$$
			$$(P - 2I)x = 0$$
			$$\left(\begin{bmatrix}4 &-1 \\ 6 & -1\end{bmatrix} - \begin{bmatrix}2 &0 \\ 0 & 2\end{bmatrix}\right)x = 0$$
			$$\begin{bmatrix}2 &-1 \\ 6 & -3\end{bmatrix}x = 0$$
			Resolviendo el sistema,
			$$\begin{bmatrix}2 &-1 \\ 6 & -3\end{bmatrix} \sim 
			\begin{bmatrix}2 &-1 \\ 0 & 0\end{bmatrix} \ra x_2 = 2x_1$$
			Luego,
			$$v = x_1v_1 + 2x_1v_2 = x_1(v_1 + 2v_2)$$
			Finalmente,
			$$B_W = \{v_1 + 2v_2\}$$
\end{enumerate}
\end{solucion}
\item Calcular los valores y vectores propios de
	$$ A = \begin{bmatrix} 1 & 2 & -1\\ 1 & 0 & 1\\ 4 & -4 & 5\end{bmatrix}$$
\begin{solucion}
Para calcular los valores propios de una matriz, debemos buscar todos los valores de $\lambda$ tales que
		$$A-\lambda I = 0$$
		Luego, para cada valor propio obtenido, debemos resolver el problema
		$$(A-\lambda I)v = 0$$
		Encontrando así los valores propios.\\
		\\
		Pasando ahora al ejercicio,
		$$\left|\begin{bmatrix}
		1-\lambda & 2 & -1\\
		1 & -\lambda & 1\\
		4 & -4 & 5-\lambda 
		\end{bmatrix}\right| = 0$$
		Usando cofactores, tenemos que
		$$\left|\begin{bmatrix}
		1-\lambda & 2 & -1\\
		1 & -\lambda & 1\\
		4 & -4 & 5-\lambda 
		\end{bmatrix}\right| = (1-\lambda)(-\lambda(5-\lambda) - -4) - 2((5-\lambda)-4) -(-4 - -4\lambda)$$
		Simplificando,
		$$=-\lambda^3 + 6\lambda^2 -11\lambda + 6 = -(\lambda - 1)(\lambda - 2) (\lambda - 3)$$
		Por lo que los valores propios son
		$$\lambda_1 =  1, \quad \lambda_2 = 2, \quad \lambda_3 = 3$$
		Busquemos ahora los vectores propios asociados a cada valor propio
		\begin{enumerate}[1)]
			\item $\lambda_1 = 1$\\
			\\
			Debemos resolver el sistema $(A-I)v = 0$, esto es
			\small$$\begin{bmatrix}
			0 & 2 & -1\\
			1 & -1 & 1\\
			4 & -4 & 4
			\end{bmatrix} \sim 
			\begin{bmatrix}
			1 & -1 & 1\\
			0 & 2 & -1\\
			4 & -4 & 4
			\end{bmatrix} \sim 
			\begin{bmatrix}
			1 & -1 & 1\\
			0 & 2 & -1\\
			0 & 0 & 0
			\end{bmatrix} \sim 
			\begin{bmatrix}
			1 & -1 & 1\\
			0 & 1 & -\frac{1}{2}\\
			0 & 0 & 0
			\end{bmatrix} \sim 
			\begin{bmatrix}
			1 & 0 & \frac{1}{2}\\
			0 & 1 & -\frac{1}{2}\\
			0 & 0 & 0
			\end{bmatrix}$$
			Por lo que
			$$\begin{array}{rcl}
			x_1 & = & -\dfrac{1}{2}x_3\\
			x_2 & = & \dfrac{1}{2}x_3\\
			x_3 & = & x_3
			\end{array} \Longrightarrow
			v_1 = \begin{pmatrix}
			-\frac{1}{2}\\
			\frac{1}{2}\\
			1
			\end{pmatrix} \rightarrow
			v_1 = \begin{pmatrix}
			-1\\
			1\\
			2
			\end{pmatrix}$$
			La solución del sistema es el generado de $v$. Por esto, lo simplificamos al final para obtener un vector propio sin fracciones.
			
			\item $\lambda_2 = 2$\\
			\\
			Debemos resolver el sistema $(A-2I)v = 0$, esto es
			\small$$\begin{bmatrix}
			-1 & 2 & -1\\
			1 & -2 & 1\\
			4 & -4 & 3
			\end{bmatrix} \sim 
			\begin{bmatrix}
			1 & 0 & \frac{1}{2}\\
			0 & 1 & -\frac{1}{4}\\
			0 & 0 & 0
			\end{bmatrix}$$
			Por lo que
			$$\begin{array}{rcl}
			x_1 & = & -\dfrac{1}{2}x_3\\
			x_2 & = & \dfrac{1}{4}x_3\\
			x_3 & = & x_3
			\end{array} \Longrightarrow
			v_2 = \begin{pmatrix}
			-\frac{1}{2}\\
			\frac{1}{4}\\
			1
			\end{pmatrix} \rightarrow
			v_2 = \begin{pmatrix}
			-2\\
			1\\
			4
			\end{pmatrix}$$
			
			\item $\lambda_3 = 3$\\
			\\
			Debemos resolver el sistema $(A-3I)v = 0$, esto es
			\small$$\begin{bmatrix}
			-2 & 2 & -1\\
			1 & -3 & 1\\
			4 & -4 & 2
			\end{bmatrix} \sim 
			\begin{bmatrix}
			1 & 0 & \frac{1}{4}\\
			0 & 1 & -\frac{1}{4}\\
			0 & 0 & 0
			\end{bmatrix}$$
			Por lo que
			$$\begin{array}{rcl}
			x_1 & = & -\dfrac{1}{4}x_3\\
			x_2 & = & \dfrac{1}{4}x_3\\
			x_3 & = & x_3
			\end{array} \Longrightarrow
			v_3 = \begin{pmatrix}
			-\frac{1}{4}\\
			\frac{1}{4}\\
			1
			\end{pmatrix} \rightarrow
			v_3 = \begin{pmatrix}
			-1\\
			1\\
			4
			\end{pmatrix}$$
		\end{enumerate}
		Finalmente, los vectores propios son
		$$v_1 = \begin{pmatrix}
		-1\\
		1\\
		2
		\end{pmatrix}, \quad
		v_2 = \begin{pmatrix}
		-2\\
		1\\
		4
		\end{pmatrix}, \quad
		v_3 = \begin{pmatrix}
		-1\\
		1\\
		4
		\end{pmatrix}$$
\end{solucion}
\item Determine si las siguientes afirmaciones son Verdaderas o Falsas.
\begin{enumerate}[a)]
\item El espacio fila de $AB$ es subespacio del espacio fila de $B$.
\item Dada una base en un espacio vectorial $V$, el vector coordenado de un vector de $V$ con respecto a esa base, es único.
\item Si $A$ y $B$ son similares y $\lambda$ es valor propio de $A$ entonces $\lambda$ es también valor propio de $B$.
\end{enumerate}
\begin{solucion}

\begin{enumerate}[a)]
\item El espacio fila de $AB$ es subespacio del espacio fila de $B$.
			
			$$C = AB  \ra C^T = B^TA^T \ra Col(C^T) = Col(B^TA^T)$$
			Como al multiplicar una matriz por otra podemos, en el mejor de los casos, obtener la misma dimensión de antes,
			$$Col(C^T) \subset Col(B^T) \ra Fil(C) \subset Fil(B)$$
			Luego,
			$$Fil(AB) \subset Fil(B)$$
			Con lo que concluimos que el espacio fila de $AB$ es subespacio del espacio fila de $B$. Entonces, la afirmación es {\bf Verdadera}.
\item Dada una base en un espacio vectorial $V$, el vector coordenado de un vector de $V$ con respecto a esa base, es único.
			
			Los vectores de una base son $L.I.$, por lo que existe una única combinación lineal para generar cada vector. Luego, la afirmación es {\bf Verdadera}.
\item Si $A$ y $B$ son similares y $\lambda$ es valor propio de $A$ entonces $\lambda$ es también valor propio de $B$.
			
			Tomemos, las matrices
			$$A = \begin{bmatrix}
			1 & 0 \\ 0 & 1
			\end{bmatrix}, \quad 
			B = \begin{bmatrix}
			2 & 0 \\ 0 & 2
			\end{bmatrix}$$
			Los valores propios de $A$ son $\lambda = 1$ (multiplicidad 2) y los de $B$ son $\lambda = 2$ (multiplicidad 2), por lo que la afirmación es {\bf Falsa}.
\end{enumerate}
\end{solucion}
\item Sea $A$ una matriz invertible, y sea $\lambda$ un valor propio de $A$. Demuestre que $\lambda \neq 0$ y que $\dfrac{1}{\lambda}$ es un valor propio de $A^{-1}$.
\begin{solucion}

		P.D. $\lambda \neq 0 \wedge \dfrac{1}{\lambda} \text{ es VP de }A^{-1}$
		
		Digamos que $\lambda = 0$. Luego, como $\lambda = 0$ es un VP, tenemos que
		$$det(A-\lambda I) = 0$$
		$$det(A-0 I) = 0$$
		$$det(A) = 0$$
		Pero $det(A) \neq 0$, ya que $A$ es invertible. Entonces, $\lambda \neq 0$.
		
		Notemos también que
		$$A^{-1}x 
		= A^{-1}\left(\dfrac{1}{\lambda}(\lambda x)\right)
		= A^{-1}\dfrac{1}{\lambda}(\lambda x)
		= \dfrac{1}{\lambda}A^{-1}(A x)
		= \dfrac{1}{\lambda} x$$
		$$A^{-1}x 
		= \dfrac{1}{\lambda} x$$
		De donde se desprende que $\dfrac{1}{\lambda}$ es valor propio de $A^{-1}$.
\end{solucion}
\item Sea $A$ una matriz de $2\times 2$ de rango 1 tal que $A\begin{bmatrix} 1 \\ 2\end{bmatrix} = \begin{bmatrix} 2 \\ 4 \end{bmatrix}$. ¿Es $A$ diagonalizable? Justifique.
\begin{solucion}
Notemos que
		$$A\begin{bmatrix} 1 \\ 2\end{bmatrix} = \begin{bmatrix} 2 \\ 4 \end{bmatrix}
		= A\begin{bmatrix} 1 \\ 2\end{bmatrix} = 2\begin{bmatrix} 1 \\ 2 \end{bmatrix}$$
		Luego, $\lambda_1 = 2$ es valor propio de la matriz $A$ y $v_1 = \begin{bmatrix} 1 \\ 2\end{bmatrix}$ es el vector propio asociado.
		
		Como $A$ es de rango 1, $\exists u \neq 0$ tal que $Au = $. Para ese $u$ se cumple que
		$$Au = 0 \ra Au = 0 u$$
		Luego, $\lambda_2 = 0$ es valor propio de la matriz $A$ y $v_2 = u$ es el vector propio asociado.
		
		Finalmente, como la multiplicidad algebraica es igual a la multiplicidad geométrica, la matriz es diagonalizable.
\end{solucion}
\item Diagonalice la matriz
	$$M = \begin{bmatrix}
	1 & 0 & 0\\
	1 & 1 & 2 \\
	1 & 0 & 3
	\end{bmatrix}$$
	y encuentre una matriz $N$ tal que $N^3 = M$.
\begin{solucion}
Buscamos los valores propios,
		$$det(M-\lambda I) = \left|\begin{bmatrix}
		1-\lambda & 0 & 0\\
		1 & 1-\lambda & 2 \\
		1 & 0 & 3-\lambda
		\end{bmatrix}\right|
		= (1-\lambda)(1-\lambda)(3-\lambda) = 0$$
		$$\lambda_1 = 1 \ra \text{multiplicidad 2}$$
		$$\lambda_2 = 3 \ra \text{multiplicidad 1}$$
		Buscamos ahora los vectores propios,
		\begin{itemize}
			\item $\lambda_1 = 1$
			
			$$(A-I)x = 0$$
			$$\begin{bmatrix}
			0 & 0 & 0\\
			1 & 0 & 2 \\
			1 & 0 & 2
			\end{bmatrix} \ra v_1 = \begin{pmatrix}
			2 \\ 0 \\ -1
			\end{pmatrix}, \quad
			v_2 = \begin{pmatrix}
			0 \\ 1 \\ 0
			\end{pmatrix}$$
			
			\item $\lambda_2 = 3$
			
			$$(A-3I)x = 0$$
			$$\begin{bmatrix}
			-2 & 0 & 0\\
			1 & -2 & 2 \\
			1 & 0 & 0
			\end{bmatrix} \ra v_3 = \begin{pmatrix}
			0 \\ 1 \\ 1
			\end{pmatrix}$$
		\end{itemize}
			Recordemos ahora que para diagonalizar una matriz debemos expresarla de la forma
			$$M = PDP^{-1} \ra P = \begin{bmatrix}
			v_1 & v_2 & v_3
			\end{bmatrix}, \quad
			D = \begin{bmatrix}
			\lambda_1 & 0 & 0 \\
			0 & \lambda_2 & 0 \\
			0 & 0 & \lambda_3
			\end{bmatrix}$$
			Es importante recordar que cada vector propio debe ir asociado con su respectivo valor propio (misma columna).
			
			En nuestro ejercicio, debemos considerar $\lambda_1 = \lambda_2 = 1$ y $\lambda_3 = 3$. Entonces,
			$$P = \begin{bmatrix}
			2 & 0 & 0 \\
			0 & 1 & 1 \\
			-1 & 0 & 1
			\end{bmatrix}, \quad
			D = \begin{bmatrix}
			1 & 0 & 0 \\
			0 & 1 & 0 \\
			0 & 0 & 3
			\end{bmatrix}, \quad
			P^{-1} = \begin{bmatrix}
			\frac{1}{2} & 0 & 0 \\
			-\frac{1}{2} & 1 & -1 \\
			\frac{1}{2} & 0 & 1
			\end{bmatrix}$$
			Sea
			$$R = \sqrt[3]{D} = \begin{bmatrix}
			1 & 0 & 0 \\
			0 & 1 & 0 \\
			0 & 0 & \sqrt[3]{3}
			\end{bmatrix}$$
			y sea
			$$N = PRP^{-1}$$
			Notemos que	
			$$N^3 = (PRP^{-1}) = PR\cancel{P^{-1}P}R\cancel{P^{-1}P}RP^{-1} = PR^3P^{-1} = PDP^{-1} = M$$
			Finalmente, $N = P\sqrt[3]{D}P^{-1}$ es la matriz buscada.
\end{solucion}
\item Diagonalice $A =
	\begin{bmatrix} 
	-1 & -2 & 2 \\
	0 & -1 & 0 \\
	0 & -2 & 1
	\end{bmatrix}$ y diagonalice $B = A^{10} + A - I$ 
\begin{solucion}
Buscamos los valores propios,
		$$det(A-\lambda I) = \left|\begin{bmatrix}
		-1-\lambda & -2 & 2 \\
		0 & -1-\lambda & 0 \\
		0 & -2 & 1-\lambda
		\end{bmatrix}\right|
		= (1+\lambda)^2(1-\lambda) = 0$$
		$$\lambda_1 = -1 \ra \text{multiplicidad 2}$$
		$$\lambda_2 = 1 \ra \text{multiplicidad 1}$$
		Buscamos ahora los vectores propios,
		\begin{itemize}
			\item $\lambda_1 = -1$
			
			$$(A+I)x = 0 \ra v_1 = \begin{pmatrix}
			1 \\ 0 \\ 1
			\end{pmatrix}$$
			
			\item $\lambda_2 = 1$
			
			$$(A-I)x = 0 \ra v_2 = \begin{pmatrix}
			1 \\ 0 \\ 0
			\end{pmatrix}, \quad
			v_2 = \begin{pmatrix}
			0 \\ 1 \\ 1
			\end{pmatrix}$$
		\end{itemize}
		Luego,
		$$A = PDP^{-1} \ra P = \begin{bmatrix}
		1 & 1 & 0 \\
		0 & 0 & 1 \\
		1 & 0 & 1
		\end{bmatrix}, \quad
		D = \begin{bmatrix}
		1 & 0 & 0 \\
		0 & -1 & 0 \\
		0 & 0 & -1
		\end{bmatrix}, \quad
		P^{-1} = \begin{bmatrix}
		0 & -1 & 1 \\
		1 & 1 & -1 \\
		0 & 1 & 0
		\end{bmatrix}$$
		Reemplacemos ahora en $B$,
		$$\begin{array}{lcl}
		B & = & A^{10} + A - I \\
		 & = & (PDP^{-1})^{10} + PDP^{-1} - I \\
		 & = & (PDP^{-1})^{10} + PDP^{-1} - PIP^{-1} \\
		B & = & P(D^10 + D - I)P^{-1}
		\end{array}$$
		Además, sabemos que
		$$D^{10} = \begin{bmatrix}
		1^{10} & 0 & 0 \\
		0 & (-1)^{10} & 0 \\
		0 & 0 & (-1)^{10}
		\end{bmatrix}
		 = \begin{bmatrix}
		1 & 0 & 0 \\
		0 & 1 & 0 \\
		0 & 0 & 1
		\end{bmatrix} = I$$
		Por lo tanto,
		$$B = P(I + D - I)P^{-1} = PDP^{-1}$$
		Con lo que concluimos que $A=B$, por lo que la diagonalización de $A$ también es diagonalización de $B$.
\end{solucion}
\item La matriz $\begin{bmatrix} 1 & -2 \\ 1 & 3 \end{bmatrix}$ actúa sobre $\mathbb{C}^2$. Determine los valores propios y una base para cada espacio propio de $\mathbb{C}^2$.
\begin{solucion}
En primer lugar, buscamos los valores propios,
		$$det(A-\lambda I) = \left|\begin{bmatrix} 
		1-\lambda & -2 \\ 
		1 & 3-\lambda \end{bmatrix}\right| = (1-\lambda)(3-\lambda) + 2 = \lambda^2 - 4\lambda + 5$$
		$$\lambda_{1,2} = 2 \pm i$$
		Como los valores propios son un par de complejos conjugados, los vectores propios asociados a cada uno también lo serán, por lo que basta con buscar un solo vector propio.
		
		Para $\lambda = 2+i$,
		$$(A-(2+i)\lambda)x = 0 \ra 
		\begin{bmatrix} 
		-1-i & -2 \\ 
		1 & 1-i 
		\end{bmatrix} \sim 
		\begin{bmatrix} 
		-1-i & -2 \\ 
		0 & 0
		\end{bmatrix}$$
		$$(-1-i)x_1 - 2x_2 = 0 \ra \begin{array}{lcl}
		x_1 & = & x_1 \\
		x_2 & = & \dfrac{-1-i}{2}x_1
		\end{array}$$
		Luego,
		$$v_1 = \begin{pmatrix}
		2 \\
		-1-i
		\end{pmatrix}, \quad
		v_2 = \begin{pmatrix}
		2 \\
		-1+i
		\end{pmatrix} $$
		Por último, las bases de los espacios propios están formadas por los vectores propios asociados a cada valor propio, es decir
		$$E_{2+i} = Gen\left\{\begin{pmatrix}
		2 \\
		-1-i
		\end{pmatrix}\right\}, \quad 
		E_{2-i} = Gen\left\{\begin{pmatrix}
		2 \\
		-1+i
		\end{pmatrix}\right\}$$
\end{solucion}
\item Sea $A$ una matriz de $5 \times 5$ tal que $A^t = -A$ y $q$ el polinomio dado por $q(x) = 2-x^2+4x^3$. Demuestre que $2$ es valor propio de la matriz $q(A)$.
\begin{solucion}
Sabemos que
		$$det(A) = det(A^T) \wedge det(A^T) = det(-A)$$
		$$det(A) = det(-A)$$
		$$det(A) = -det(A)$$
		$$det(A) = 0$$
		Luego, existe $u \neq 0 | Au = 0$
		
		Notemos que
		$$\begin{array}{lcl}
		q(A) & = & 2I - A^2 + 4A^3\\
		q(A) u & = & (2I - A^2 + 4A^3)u\\
		& = & 2Iu - A^2u + 4A^3u \\
		& = & 2Iu - A(Au) + 4A^2(Au) \\
		& = & 2Iu - A0 + 4A^20 \\
		& = & 2Iu \\
		q(A) u & = & 2u
		\end{array}$$
		Luego, $\lambda = 2 $ es $VP$ de $q(A)$
\end{solucion}
\end{preguntas}
\end{document}